\begin{helper}
Let \(f:\mathbb R\to\mathbb R\), and let \(a\le b\).  Assume that \(f\) is
interval-integrable on \([a,b]\) and that \(f(x)\ge0\) for every
\(x\in(a,b]\).  Then
\[
\int_{(a,b]}^{+}\operatorname{ofReal}(f(x))\,dx
=
\operatorname{ofReal}\left(\int_a^b f(x)\,dx\right).
\]
\end{helper}
\begin{proof}
Because \(a\le b\), the interval integral \(\int_a^b f(x)\,dx\) is the
Bochner integral of \(f\) over the restricted Lebesgue measure on
\((a,b]\).  The endpoint \(a\) has measure zero, so using \((a,b]\) instead
of \([a,b]\) does not change either side.  The interval-integrability
hypothesis says precisely that \(f\) is integrable for this restricted finite
measure.  The pointwise nonnegativity hypothesis gives \(0\le f\) almost
everywhere for the same restricted measure.  The standard
Lebesgue-integral conversion theorem for an integrable nonnegative real-valued
function states that the nonnegative integral of
\(\operatorname{ofReal}\circ f\) equals \(\operatorname{ofReal}\) of the real
integral of \(f\).  Applying that theorem to the restricted measure
\(\operatorname{vol}|_{(a,b]}\) gives exactly the displayed equality.
\end{proof}
